% TOP_PROBLEM: {"id":30007574,"problem_number":"LOCAL-30007574","title":"Mechanisms with resale","category_id":21,"category":{"id":21,"name":"economics","display_name":"Economics","description":"Open economic theory statements, published research agendas, and proposed scoped specifications; question provenance and historical priority are qualified per record.","slug":"economics","order_index":21,"created_at":"2026-10-08T00:00:00Z"},"status":"research_question","external_url":null,"legacy_ids":[],"published":false,"statement":"Jointly design the initial allocation mechanism and information disclosure when buyers subsequently trade in a specified resale game, accounting for total incentives and posterior beliefs.","economics":{"original_id":63,"field":"Mechanism design and market design","kind":"design","formalization_basis":"adapted_research_question","source_status":"research_agenda","priority_status":"earliest_found","priority_note":"This is the earliest source in which the relevant agenda was verified during this audit; absolute priority has not been established.","earliest_verified_reference":{"authors":["Dirk Bergemann","Juuso Välimäki"],"title":"Dynamic Mechanism Design: An Introduction","year":2018,"url":"https://cowles.yale.edu/sites/default/files/2022-09/d2102-r.pdf","doi":"10.1257/jel.20180892","locator":"§8 Concluding Remarks, pp.49–51 of June19,2018 revision","evidence_summary":"The conclusion explicitly identifies weaker common-prior assumptions, sequential participation, outside resale, and limited adoption of long-term contracts as research questions."},"current_reference":{"authors":["Dirk Bergemann","Juuso Välimäki"],"title":"Dynamic Mechanism Design: An Introduction","year":2018,"url":"https://cowles.yale.edu/sites/default/files/2022-09/d2102-r.pdf","doi":"10.1257/jel.20180892","locator":"§8 Concluding Remarks, pp.49–51 of June19,2018 revision","evidence_summary":"The conclusion explicitly identifies weaker common-prior assumptions, sequential participation, outside resale, and limited adoption of long-term contracts as research questions."},"formalization_note":"The fixed resale game and joint disclosure optimization are specified here. Only the broader information–mechanism interaction agenda is verified in §8."},"statusQualification":"Published question or research agenda verified at the cited locator. The mathematical specification is adapted for this catalogue and includes additional assumptions; source publication does not establish first-ever priority or certify this exact model remains unsolved. The fixed resale game and joint disclosure optimization are specified here. Only the broader information–mechanism interaction agenda is verified in §8.","statusReviewedAt":"2026-10-08"}
% FIRST_POSTED: 2026-10-08
% ENTRY_KIND: statement_only
% !TeX program = lualatex
\documentclass[11pt]{article}
\usepackage{fontspec}
\usepackage[a4paper,margin=25mm]{geometry}
\usepackage{amsmath,amssymb,mathtools}
\usepackage{xurl}
\usepackage[hidelinks]{hyperref}
\newcommand{\sref}[2]{\hyperlink{source:#1:#2}{[S#2]}}
\setlength{\emergencystretch}{3em}
\begin{document}
% problemId:30007574
\section*{Mechanisms with resale}\label{30007574}
\subsection{Definitions and mathematical statement}
A seller initially owns one indivisible good. There are \(n\geq2\) buyers with private persistent values \(v_i\) drawn from a specified joint distribution \(F\) on \([0,1]^n\). The seller can commit to an initial direct allocation rule \(x(v)\), payments \(p(v)\), and a disclosure rule \(d\) that produces public or private signals about the reports. After the initial allocation, the current owner may trade the good in a fixed resale protocol \(R\), such as a finite sequential bargaining game. Specify timing, bargaining power, available information, and an equilibrium-selection rule for \(R\). Let \(U_i(v_i,\widehat v_i;M,d,R)\) be buyer \(i\)'s expected total utility from truthful value \(v_i\), initial report \(\widehat v_i\), initial transfers, and the resale equilibrium. The target is
\[
 \sup_{M,d}E_F\left[\sum_i p_i(v)\right]
 \quad\text{subject to}\quad
 U_i(v_i,v_i;M,d,R)\geq U_i(v_i,\widehat v_i;M,d,R)
\]
for every buyer and report, together with initial participation and allocation feasibility. Characterize how optimal disclosure and allocation depend on the resale protocol and value correlation, and identify when resale weakens or strengthens initial-market revenue and efficiency. Disclose all equilibrium assumptions because posterior beliefs and trading incentives depend on information released by the mechanism. This is a research-agenda adaptation about interactions with outside markets, not a claim that auctions with resale have no established characterization in specific environments.

\par\textbf{Formulation and scope.} The fixed resale game and joint disclosure optimization are specified here. Only the broader information–mechanism interaction agenda is verified in §8.

\par\textbf{Open-question provenance.} An explicit question or research agenda was verified at the source locator. This does not by itself certify that every later version remains unresolved \sref{30007574}{1}.

\par\textbf{Historical priority.} This is the earliest source in which the relevant agenda was verified during this audit; absolute priority has not been established.

\subsection{Short English statement}
Jointly design the initial allocation mechanism and information disclosure when buyers subsequently trade in a specified resale game, accounting for total incentives and posterior beliefs.

\subsection{Sources}
\begin{itemize}\raggedright
\item[\textnormal{[S1]}]\hypertarget{source:30007574:1}{} Dirk Bergemann, Juuso Välimäki. 2018. Dynamic Mechanism Design: An Introduction. §8 Concluding Remarks, pp.49–51 of June19,2018 revision.
\url{https://cowles.yale.edu/sites/default/files/2022-09/d2102-r.pdf}.
DOI: 10.1257/jel.20180892.
{\small\textit{Earliest verified question/agenda reference; historical priority qualified below. The conclusion explicitly identifies weaker common-prior assumptions, sequential participation, outside resale, and limited adoption of long-term contracts as research questions.}}
\end{itemize}
\end{document}
